\documentclass[10pt]{article}
\usepackage[margin=20mm]{geometry}
\usepackage{amsmath,amssymb,mathtools,fontspec,unicode-math,xurl,hyperref}
\IfFontExistsTF{DejaVu Serif}{\setmainfont{DejaVu Serif}}{\setmainfont{Latin Modern Roman}}
\setmathfont{latinmodern-math.otf}
\setlength{\parindent}{0pt}\setlength{\parskip}{6pt}\sloppy
\title{Grading and proof audit of
Exotic Dehn twists and homotopy coherent group actions}
\author{}
\date{}
\begin{document}
\maketitle
Abstract: Two explicit suspension calculations expose an inconsistent normalization in the family Frøyshov invariant used in the proof of the main theorem. The printed definition changes by two under a trivial real two-dimensional suspension, whereas Lemma 5.15 predicts no change. A separate complex-line test changes the proposed boundary invariant by four while preserving its stable presentation. These establish one Major defect in the core derivation as printed; they do not refute the asserted infinite-order Dehn twist. The higher-coherence step of Lemma 3.3 remains a separate Needs Evidence item.\par
\section*{Identity, version, and scope}
The reviewed paper is by Sungkyung Kang, JungHwan Park, and Masaki Taniguchi, DOI 10.1007/s00222-025-01378-1 (https://doi.org/10.1007/s00222-025-01378-1) [KPT]. The source is the supplied 33-page publisher PDF, first published online on 10 October 2025. Throughout, PDF p. \(r\) means the \(r\)-th page of that exact file; its printed page is \(208+r\). The frozen identifier is M25-0114.\par
All 33 pages of the extracted source text were read in this continuation. The detailed checks concern theorem hypotheses, the collar construction and finite-order reduction, the definitions and suspension behavior in Sections 5.2–5.5, and their use in the final contradiction. Reading the analytic construction in Section 4 is recorded as a dependency review, not a fresh proof of its Fredholm, compactness, or Conley-index results. The external gauge-theoretic and lattice theorems cited by the authors were not independently reproved.\par
Theorem 1.1 (PDF p. 2, printed p. 210) assumes a smooth compact positive-definite filling \(X\) of the oriented Brieskorn homology sphere \(Y\ne S^3\), with \(b_1(X)=0\). Its conclusion concerns infinite order in \(\pi_0\operatorname{Diff}^+(X,\partial X)\). Simple connectivity is an additional hypothesis of the exoticness conclusion in Corollary 1.2 (PDF p. 3), not a substitute for the weaker hypothesis in Theorem 1.1.\par
\section*{Problem 1: incompatible grading conventions in a core invariant}
\subsection*{Source statements and exact locations}
Equation reference: eq:raw. Definition 5.3 and Definition 5.5 occur on PDF p. 24 (printed p. 232), with the Thom-class explanation continuing on p. 25. Write the quantity in the displayed Definition 5.5 as \[
 d^{\mathrm{raw}}_{j,n}(X)
 =\min\{i:\text{an admissible class lies in }H^i_{S^1}(X,s(B))\}-2j .\] Admissibility means that restriction to the fixed-point sphere bundle has the stated leading term \(S^jU^k\) times its Thom generator, modulo the displayed ideal. There is no subtraction of the fixed-point level in [eq:raw] (\#eq:raw). The equivariant definition on PDF p. 23 likewise uses an ordinary cohomological degree.\par
Equation reference: eq:printedshift. For \(E=E_{\mathbb R}\oplus E_{\mathbb C}\), Lemma 5.15 (PDF p. 27, printed p. 235) asserts \[
 d_{j,n}(\Sigma^E X)
 =\dim_{\mathbb R}E_{\mathbb C}+d_{j,n}(X).\] Here \(S^1\) acts trivially on \(E_{\mathbb R}\) and by scalar multiplication on \(E_{\mathbb C}\). Its proof displays a Thom isomorphism shifting total degree by \(\dim_{\mathbb R}E\), which does not by itself yield [eq:printedshift] (\#eq:printedshift) under [eq:raw] (\#eq:raw).\par
\subsection*{Check A: a fully specified counterexample to the real-suspension formula}
Equation reference: eq:example. Fix \[p=3,\qquad n=2,\qquad B=(B\mathbb Z_3)_2,\qquad j=0,\] and take \[
 X=B\times S^0,\qquad E_{\mathbb R}=B\times\mathbb R^2,\qquad E_{\mathbb C}=0.\] Use the point at infinity as section and give \(X\) and \(E_{\mathbb R}\) the trivial circle action.\par
The two-skeleton used by the paper has finitely many cells, so \(B\) is a compact connected finite CW complex. The projection of \(X\) is a fiber bundle with finite CW fiber. Its fixed-point bundle is \(B\times S^0\), the compactification of the rank-zero bundle. The complement of the fixed points is empty, and hence satisfies the required freeness condition. The bundle \(E_{\mathbb R}\) is oriented. The suspended family \(\Sigma^E X=B\times S^2\), with its infinity section, satisfies the same conditions with fixed-point level two. Thus every assumption used in Definition 5.3 and Lemma 5.15 holds.\par
Equation reference: eq:value0, eq:value2. Set \(A=H^*(B\times BS^1;\mathbb F_3)\). For the first family, \[H^*_{S^1}(X,s(B);\mathbb F_3)\cong A.\] The fixed-point restriction is the identity. The degree-zero unit is admissible for \(j=k=0\); there are no negative-degree classes. Consequently, \[
 d^{\mathrm{raw}}_{0,2}(X)=0.\] For the suspended family the Thom isomorphism gives \[H^*_{S^1}(\Sigma^E X,s(B);\mathbb F_3)\cong A[-2].\] The Thom generator has degree two and unit coefficient, and restriction is again the identity. Therefore \[
 d^{\mathrm{raw}}_{0,2}(\Sigma^E X)=2.\] In contrast, the right-hand side of [eq:printedshift] (\#eq:printedshift) is \(0+0=0\). Equations [eq:value0] (\#eq:value0)–[eq:value2] (\#eq:value2) are an analytic counterexample to the lemma with the printed degree convention. The argument uses \(j=0\), so it is unaffected by truncation of high powers of the base class \(S\). The same degree calculation works over either field appearing in the surrounding printed notation; the field label itself is not being counted as a substantive finding.\par
\subsection*{Check B: the boundary normalization fails a stable-presentation test}
Equation reference: eq:Dplus. On PDF p. 13 (printed p. 221), repeated on p. 22, the stable Floer presentation is \[SWF(Y,t,g_Y)=\Sigma^{-V_\lambda^0(g_Y)}I_\lambda^\mu(g_Y).\] The definition at the top of PDF p. 30 (printed p. 238) uses a plus sign: \[
 D(I,V)=d_{j,n}(I)-2c+\dim_{\mathbb R}V,
 \qquad c=n(Y,t,g_Y).\] It is immediately followed by a claim of independence of the spectral cutoffs and by Corollary 5.19 comparing it with the equivariant invariant.\par
Equation reference: eq:four. There is a direct formal compatibility test. Take \(j=0\), keep \(c\) fixed, and let \(W=\mathbb C\) have circle weight one and trivial \(\mathbb Z_3\)-action. Compare the presentations \[(I,V)=(S^0,0),\qquad (I',V')=(S^{\mathbb C},\mathbb C),\] with the trivial families over \(B=(B\mathbb Z_3)_2\). Both represent the same stable object: \[\Sigma^{-\mathbb C}S^\mathbb C\simeq S^0.\] The equivariant Thom generator of \(S^\mathbb C\) has degree two, and its fixed-point restriction is its Euler class \(U\). Thus \(d_{0,2}(I')=2\), whereas \(d_{0,2}(I)=0\). Both fixed-point levels are zero, so subtracting the fixed-point level would not change this test. Formula [eq:Dplus] (\#eq:Dplus) gives \[
 D(I',V')-D(I,V)=2+2=4.\] The sign is therefore incompatible with invariance of the stable presentation. Interpreting the dimension term as complex dimension would give a change of three, not zero.\par
This calculation is a formal suspension test of the displayed normalization. It is not a numerical reconstruction of a particular three-manifold’s spectral cutoffs. No actual gauge-theoretic spectrum has been computed or attributed to the authors.\par
\subsection*{Alternative readings, local repairs, and their limits}
An intended level-normalized convention \[\bar d_{j,n}(X)=d^{\mathrm{raw}}_{j,n}(X)
                -\operatorname{level}(X)\] removes the discrepancy in Check A. Under the printed raw convention, the trivial real two-dimensional suspension instead has to contribute two. Both statements follow from the explicit Thom calculation. Neither convention is supplied by the displayed Definition 5.5.\par
For Check B, replacing the plus sign by a minus sign cancels the tested complex-line shift: \[[d_{0,2}(I')-\dim_{\mathbb R}V']-
 [d_{0,2}(I)-\dim_{\mathbb R}V]=2-2=0.\] This is a verified local repair for that test only. A coordinated repair still has to state one grading convention, reconcile the equivariant and family definitions, and carry it through the Bauer–Furuta map and the comparison theorem. This report does not assert that those global repairs have been completed.\par
There is related degree bookkeeping in Proposition 5.18 (PDF pp. 28–29): the proof sets \(k=d_j(X)\) and places the defining class in degree \(k\), although [eq:raw] (\#eq:raw) puts it in degree \(d_j(X)+2j\). For fixed \(j\), choosing a skeleton above the actual class degree addresses that local mismatch. This observation is part of the same grading finding, not a separate claim that eventual comparison is false.\par
\subsection*{Core dependence and severity}
Theorem 5.20 on PDF p. 30 explicitly invokes Lemmas 5.15 and 5.17 and then Corollary 5.19 to establish the family Frøyshov inequality. Lemma 6.1 on PDF p. 31 uses that inequality, and the concluding contradiction on PDF p. 32 uses Lemma 6.1 to prove Theorem 1.1. Thus the faulty normalization is directly inside the core proof chain. The supported judgment is Major: an invalid core derivation under the definitions as printed. The evidence does not establish falsity of the infinite-order or exoticness conclusions.\par
\section*{Problem 2: higher coherence requires additional evidence}
Lemma 3.3 (PDF p. 11, printed p. 219) constructs the two-simplex homotopies from a collar diffeomorphism \(\tau\) and a path \(H_t\) satisfying \[\tau^p=\Phi,\qquad H_0=\Phi,\qquad H_1=\operatorname{id},\qquad
 H_t\tau=\tau H_t.\] The displayed three-simplex comparison follows by strict commutation. The passage to all higher simplices is attributed to contractibility of the interval-reparametrization space.\par
The stated algebraic relations alone do not guarantee that extension for arbitrary topological groups. To see the distinction, take a topological-group model \(G\simeq\Omega S^2\), put \(\tau=1\), and choose \(H\) as a loop generating \(\pi_1(G)\cong\mathbb Z\). The cyclic carry rule for the two-simplex homotopies represents the generator \(t\in H^2(B\mathbb Z_p;\mathbb Z)=\mathbb Z/p\). If that assignment extended to \(f:B\mathbb Z_p\to BG\simeq S^2\), it would pull the generator \(u\in H^2(S^2;\mathbb Z)\) back to \(t\). But \[0=f^*(u^2)=t^2,\] whereas \(t^2\ne0\) in \(H^4(B\mathbb Z_p;\mathbb Z)=\mathbb Z/p\). Thus a general inference from the commuting relation to all higher fillers is insufficient.\par
This diagnostic is deliberately outside the paper’s manifold class: it is not a four-manifold counterexample. The particular collar construction may provide additional geometry that resolves the obstruction. The unresolved item is an explicit compatible higher extension for those collar data, or a theorem establishing it. The judgment is Needs Evidence, with no confirmed-error severity assigned. Its dependency is Lemma 3.3 to Proposition 3.5 and then to the final finite-order contradiction.\par
\section*{Checks that do not support additional defects}
\subsection*{Finite order, prime choice, and monodromy}
The collar calculation in Lemma 3.1 commutes strictly: the first isotopy and \(\tau\) are circle rotations on each collar fiber; the second isotopy has support disjoint from that of \(\tau\). For a hypothetical finite positive order \(d\), choose an odd prime \(p\) large enough for the cited Floer calculation and coprime both to \(d\) and to the Seifert multiplicity product. Choose \(d'=kd\) with \(d'\equiv1\pmod p\). Reparametrizing the circle action by \(d'\) leaves its \(\mathbb Z_p\) generator unchanged and replaces the boundary twist by \(\Phi^{d'}\), which is relatively isotopic to the identity. This verifies the reduction in Remark 3.6 (PDF p. 12). The ordinary isotopy of \(\tau\) to the identity is not required to fix the boundary, so the monodromy conclusion of Lemma 3.4 is consistent.\par
\subsection*{Cyclic Spin\(^c\) lifting}
The cohomology-isomorphism hypothesis printed in Corollary 2.5 (PDF pp. 8–9) is not an available hypothesis for the main fillings. However, the needed cyclic lifting can be obtained directly from Lemma 2.4. With \(b_1(X)=0\), the gauge group is homotopy equivalent to the constant \(S^1\) gauges. These constants are central, so the obstruction has untwisted coefficients and lies in \(H^3(B\mathbb Z_p;\mathbb Z)=0\). The lift torsor is \(H^2(B\mathbb Z_p;\mathbb Z)\), and restriction to the homology-sphere boundary induces the identity on the constant gauges. One can therefore adjust the interior lift to the chosen boundary lift. This supplies a direct explanation for the cyclic application; the citation mismatch is not treated as a failed main conclusion. The broader wording of Theorem 5.20, which does not repeat every Section 4 assumption, has not been separately disproved.\par
\subsection*{Orientation and finite skeletons}
A positive-definite filling enters the negative-definite cobordism argument with reversed orientation. With that convention, the sign of the characteristic square in Lemma 6.1 is consistent. The form on the free quotient is unimodular for a homology-sphere boundary; finite torsion does not itself invalidate the lattice step. The external characteristic-vector bound and equivariant Floer values remain cited dependencies.\par
Proposition 5.18 requires a skeleton threshold for each fixed \(j\). The final argument does not need one threshold uniform in \(j\). Likewise, the obstruction prime may be chosen after the hypothetical finite order is specified. Neither quantifier was found to create an additional contradiction.\par
\section*{Version checks, conclusion, and remaining work}
The current publisher landing page and a targeted title search were consulted for correction notices. No correction notice was located in those checks. The arXiv record inspected lists v2, dated 26 September 2024 [KPTarxiv]; it is earlier than the supplied publisher PDF. No unexamined author version is being used to certify a repair.\par
Item | Judgment | Repair/adoption status\par
Family grading | Major, demonstrated | Local test-case changes verified; global reconciliation unverified.\par
Higher coherence | Needs Evidence | Compatible higher fillers for the specific collar construction remain unverified.\par
Prime and scope checks | No issue found in checked steps | No change proposed.\par
Scientific acceptance | Pending | No human review or author adoption has been recorded.\par
The remaining scientific work is to reconcile the degree normalization through Theorem 5.20 and resolve the specific higher-coherence construction. A later independent focused acceptance is still pending.\par
\section*{Supporting checks}
67a0ef586af0fbd27edfd3f335450e36296bfd77cc34a29669f8677736cda5ad.\par
The report uses preserved source bytes and existing page renders. The current continuation supplies the scientific checks and report. The prior partial extraction is not represented as a prior scientific review.\par
The primary session read all source text and inspected the source evidence pages recorded in `source\_inspection.json`. Their roles and actual inspection scopes are recorded separately. No author code, external model review, numerical simulation, or human acceptance was performed.\par
9 S. Kang, J. Park, and M. Taniguchi, Exotic Dehn twists and homotopy coherent group actions, Inventiones mathematicae 243, 209–241. Publisher PDF first published online 10 October 2025. https://doi.org/10.1007/s00222-025-01378-1. S. Kang, J. Park, and M. Taniguchi, Exotic Dehn twists and homotopy coherent group actions, arXiv:2409.11806v2, 26 September 2024. Version-history metadata consulted; this report’s mathematical findings are bound to the publisher PDF. https://arxiv.org/abs/2409.11806.\par
\end{document}
